%-------------------------------------------------------------------------------
%   CHAPTER 4 - RESULTS
%-------------------------------------------------------------------------------

\section{Судалгааны үр дүн}

\subsection{Системийн хэрэгжүүлэлт}
"Tdelivery" системийг амжилттай хөгжүүлж, үндсэн функцуудыг бүрэн ажиллагаанд оруулсан. Систем нь хэрэглэгчдэд ойлгомжтой, ашиглахад хялбар интерфэйстэй бөгөөд дараах үндсэн хэсгүүдээс бүрдэнэ.

\subsubsection{Хэрэглэгчийн интерфэйс (UI)}
\begin{itemize}
    \item \textbf{Нэвтрэх хэсэг:} Утасны дугаараар баталгаажуулах (OTP) систем нь хэрэглэгчийн аюулгүй байдлыг хангаж, бүртгэлийн процессыг хурдасгасан. Мөн бүртгэлийн явцад гарч болох алдааг илрүүлэх зорилгоор `FirebaseDebugScreen` хэрэгслийг хөгжүүлж, системийн холболт болон эрхийн тохиргоог шалгах боломжийг бүрдүүлсэн.
    \item \textbf{Нүүр хуудас:} Бараа бүтээгдэхүүнийг ангиллаар харах, хайлт хийх, онцлох барааг харах боломжтой.
    \item \textbf{Сагс болон Захиалга:} Сагсанд бараа нэмэх, тоо хэмжээг өөрчлөх, захиалга баталгаажуулах үйл явц нь 3-аас доош алхмаар хийгддэг.
    \item \textbf{Хяналтын самбар:} Захиалагч, нийлүүлэгч, хүргэгч тус бүр өөрт хамааралтай мэдээллийг (захиалгын төлөв, орлого, хүргэлтийн түүх) харах боломжтой.
\end{itemize}

\subsection{Хоёр талт захиалгын удирдлага (Two-Sided Order Management)}
Системийн хамгийн чухал хэсэг болох захиалгын удирдлагыг "Захиалагч" болон "Хүргэгч" гэсэн хоёр талын оролцоотойгоор бүрэн шийдсэн.

\subsubsection{Захиалагчийн хэсэг (My Orders)}
Захиалагч нь "Миний захиалгууд" цэсээр дамжуулан захиалгын төлөвийг гурван үе шаттайгаар хянах боломжтой:
\begin{enumerate}
    \item \textbf{Хүлээгдэж буй (Pending):} Шинээр үүсгэсэн, хүргэгч хараахан аваагүй захиалгууд. Шар өнгийн төлөвтэй харагдана.
    \item \textbf{Хүргэлтэнд (In Delivery):} Хүргэгч хүлээн авсан захиалгууд. Энэ үед хүргэгчийн зураг, нэр, утасны дугаар харагдах бөгөөд "Залгах" товчлуур идэвхжинэ. Улбар шар эсвэл цэнхэр өнгөөр ялгарна.
    \item \textbf{Хүргэгдсэн (Delivered):} Амжилттай хүргэгдсэн захиалгууд. Ногоон өнгөөр тэмдэглэгдэнэ.
\end{enumerate}

\subsubsection{Хүргэгчийн хэсэг (My Deliveries)}
Хүргэгч нь "Миний хүргэлтүүд" цэсээр дамжуулан өөрийн хариуцаж буй ажлуудыг удирдана:
\begin{itemize}
    \item \textbf{Идэвхтэй хүргэлт:} "Шинэ" төлөвтэй захиалгыг "Хүлээн авах" товчийг дарснаар тухайн захиалга хүргэгчийн идэвхтэй жагсаалт руу шилжинэ. Энэ үед захиалагчийн хаяг, утасны дугаар харагдана.
    \item \textbf{Хүргэлт дуусгах:} Барааг хүлээлгэн өгсний дараа "Ирсэн" (Mark as Delivered) товчийг дарснаар захиалгын төлөв "Delivered" болж өөрчлөгдөнө.
\end{itemize}

\subsection{Кодчиллын хэсэг}
Системийн гол функцуудыг хэрэгжүүлсэн кодын хэсгээс тайлбарлая.

\subsubsection{Firebase тохиргоо ба Холболт}
Систем нь Firebase-тай холбогдохдоо дараах тохиргоог ашигладаг. Энэ нь аюулгүй байдлын үүднээс орчны хувьсагч (Environment Variables) ашиглан нууцлагдсан байдаг.

\begin{lstlisting}[language=JavaScript, caption=Firebase тохиргоо]
// src/config/firebase.ts
import { initializeApp } from 'firebase/app';
import { getFirestore } from 'firebase/firestore';
import { getAuth } from 'firebase/auth';

const firebaseConfig = {
  apiKey: process.env.EXPO_PUBLIC_FIREBASE_API_KEY,
  authDomain: process.env.EXPO_PUBLIC_FIREBASE_AUTH_DOMAIN,
  projectId: process.env.EXPO_PUBLIC_FIREBASE_PROJECT_ID,
  storageBucket: process.env.EXPO_PUBLIC_FIREBASE_STORAGE_BUCKET,
  messagingSenderId: process.env.EXPO_PUBLIC_FIREBASE_MESSAGING_SENDER_ID,
  appId: process.env.EXPO_PUBLIC_FIREBASE_APP_ID,
};

const app = initializeApp(firebaseConfig);
export const db = getFirestore(app);
export const auth = getAuth(app);
\end{lstlisting}

\subsubsection{Захиалга үүсгэх функц}
Захиалга үүсгэх үед өгөгдлийн бүрэн бүтэн байдлыг хангахын тулд Firestore Transaction ашигласан. Энэ нь барааны үлдэгдэл шалгах, хасах, захиалга үүсгэх үйлдлүүдийг нэгэн зэрэг гүйцэтгэдэг.

\begin{lstlisting}[language=JavaScript, caption=Захиалга үүсгэх Transaction]
// src/services/orderService.ts
export const createOrder = async (orderData) => {
  try {
    await runTransaction(db, async (transaction) => {
      // 1. Check product stock
      for (const item of orderData.items) {
        const productRef = doc(db, 'products', item.productId);
        const productDoc = await transaction.get(productRef);
        
        if (!productDoc.exists()) {
          throw "Product not found!";
        }
        
        const newStock = productDoc.data().stock - item.quantity;
        if (newStock < 0) {
          throw "Insufficient stock!";
        }
        
        // 2. Deduct stock
        transaction.update(productRef, { stock: newStock });
      }

      // 3. Create order
      const newOrderRef = doc(collection(db, 'orders'));
      transaction.set(newOrderRef, {
        ...orderData,
        createdAt: serverTimestamp(),
        status: 'pending'
      });
    });
    return { success: true };
  } catch (e) {
    console.error("Transaction failed: ", e);
    return { success: false, error: e };
  }
};
\end{lstlisting}

\subsubsection{Бодит цагийн хяналт (Real-time Listener)}
Захиалгын жагсаалтыг бодит цагт шинэчлэхийн тулд \texttt{onSnapshot} функцийг ашигласан. Энэ нь сервер дээр өөрчлөлт орох бүрт клиентийг автоматаар шинэчилдэг.

\begin{lstlisting}[language=JavaScript, caption=Бодит цагийн захиалгын жагсаалт]
// src/screens/OrdersScreen.tsx
useEffect(() => {
  const q = query(
    collection(db, 'orders'),
    where('customerId', '==', user.uid),
    orderBy('createdAt', 'desc')
  );

  const unsubscribe = onSnapshot(q, (snapshot) => {
    const ordersList = snapshot.docs.map(doc => ({
      id: doc.id,
      ...doc.data()
    }));
    setOrders(ordersList);
  });

  return () => unsubscribe(); // Cleanup function
}, [user.uid]);
\end{lstlisting}

\subsection{Процессын цахимжуулалтын үр дүн}
Уламжлалт гар ажиллагаатай процессыг цахимжуулснаар дараах үр дүнгүүд гарсан.

\begin{table}[h]
\centering
\caption{Уламжлалт болон Цахим системийн харьцуулалт}
\begin{tabular}{|p{4cm}|p{5cm}|p{5cm}|}
\hline
\textbf{Үзүүлэлт} & \textbf{Уламжлалт арга} & \textbf{"Tdelivery" систем} \\
\hline
Захиалга авах & Утас, мессежээр (3-5 минут) & Аппликейшнээр (30 секунд) \\
\hline
Мэдээллийн нарийвчлал & Алдаа гарах магадлал өндөр & Алдаа гарахгүй (Систем хянана) \\
\hline
Захиалгын хяналт & Боломжгүй эсвэл хүндрэлтэй & Бодит цагт хянах боломжтой \\
\hline
Тайлан тооцоо & Гараар нэгтгэдэг (Цаг их ордог) & Автомат тайлан (Шууд харах) \\
\hline
\end{tabular}
\end{table}

\subsection{Цахим баримт бичиг (e-POD) нэвтрүүлэлт}
Цахим баримт бичгийн системийг нэвтрүүлснээр цаасан баримтын хэрэглээг халж, хүргэлтийн баталгаажуулалтыг найдвартай болгосон.
\begin{itemize}
    \item \textbf{Зурган баталгаажуулалт:} Хүргэгч барааг хүлээлгэн өгөхдөө зураг авч системд оруулснаар маргаантай асуудлуудыг 90\% бууруулсан.
    \item \textbf{Гарын үсэг:} Хүлээн авагч дижитал гарын үсэг зурж баталгаажуулах боломжтой.
    \item \textbf{Байршлын мэдээлэл:} Хүргэлт хийгдсэн цэгийн GPS координатыг хадгалснаар хүргэлтийн үнэн зөв байдлыг баталгаажуулна.
\end{itemize}

\subsection{Маршрут оновчлолын үр дүн}
Маршрут оновчлолын алгоритмыг туршсанаар хүргэлтийн нийт хугацаа болон тээврийн зардлыг бууруулах боломжтойг харуулсан.
\begin{itemize}
    \item \textbf{Хүргэлтийн хугацаа:} Оновчтой маршрут ашигласнаар дунджаар 15-20\% хэмнэсэн.
    \item \textbf{Шатахууны зардал:} Зам тээврийн оновчтой төлөвлөлт нь шатахууны зардлыг 10-15\% бууруулах тооцоолол гарсан.
\end{itemize}

\subsection{Хэлэлцүүлэг}
Системийг хөгжүүлэх явцад гарсан гол бэрхшээлүүд болон шийдлүүд:
\begin{itemize}
    \item \textbf{Бэрхшээл:} Интернэтгүй орчинд ажиллах (Offline mode).
    \item \textbf{Шийдэл:} Firebase-ийн offline persistence боломжийг ашиглан сүлжээгүй үед өгөгдлийг кэшлэх, сүлжээнд холбогдмогц синхрончлох шийдлийг нэвтрүүлсэн.
\end{itemize}

\vfill